% GENERATED FILE. Edit canonical lessons, then run npm run generate:books.
% canonical-source-hash: fnv1a64:6537e1fa9c711eab
% canonical-lessons: BN-C436-ongkur, BN-C436-mukul, BN-C436-phota, BN-C436-gojano, BN-C436-pota

\chapter{Buds and Sprouts}
\label{ch:bn-buds-and-sprouts}
\hlchaptermodality{436}

\begin{chapteropening}
\textbf{By the end of this chapter:} \emph{I can say seeds sprouted, trees came into bud and flowers bloomed, and plant a tree.}

Complete the last lesson of chapter 436: I can say seeds sprouted, trees came into bud and flowers bloomed, and plant a tree.
\end{chapteropening}

\section[\texorpdfstring{\bn{অঙ্কুর} (ôngkur)}{অঙ্কুর (ôngkur)}]{\bn{অঙ্কুর} (ôngkur) — a sprout; a shoot}

\label{lesson:BN-C436-ongkur}



\begin{quote}
Before the new one: say the Bengali for wood, then the Bengali for to trim.
\end{quote}

\subsection*{You'll want to know: \bn{অঙ্কুর}}
\textbf{\bn{অঙ্কুর}} — \emph{ôngkur} — \textquotedblleft{}a sprout; a shoot\textquotedblright{}.

\begin{quote}
\bn{বৃষ্টির} \bn{পরে} \bn{বীজ} \bn{থেকে} \bn{অঙ্কুর} \bn{বেরোল।}
\end{quote}

\begin{quote}
\emph{brishṭir pôre bīj theke ôngkur berolô.}
\end{quote}

\begin{quote}
After the rain shoots came up from the seeds.
\end{quote}

\begin{cousinweb}
\textbf{\bn{অঙ্কুর}} (\emph{ôngkur}) is Sanskrit \emph{aṅkurá}, \textquotedblleft{}a sprout, a shoot\textquotedblright{} (CDIAL 109), taken over whole; CDIAL also lists the inherited Bengali \emph{ã̄kur}.
\end{cousinweb}

\begin{grammarlens}[title={What you've built}]
One more word for getting around.
\end{grammarlens}

\subsection*{Your turn}
\begin{itemize}
\raggedright
  \item \emph{Say it:} \emph{ôngkur}
  \item \emph{Say it:} \emph{ôngkur}, once more
  \item \emph{Say it:} \emph{brishṭir pôre bīj theke ôngkur berolô.} — After the rain shoots came up from the seeds.
  \item \emph{Recall:} say the Bengali for wood, then the Bengali for to trim, then say \emph{ôngkur} again
\end{itemize}

\begin{tcolorbox}[breakable,colback=teal!4,colframe=teal!35!black,title={Before you move on}]
What does \bn{অঙ্কুর} mean? (\textquotedblleft{}A sprout; a shoot\textquotedblright{}.) Say it once more.
\end{tcolorbox}
\section[\texorpdfstring{\bn{মুকুল} (mukul)}{মুকুল (mukul)}]{\bn{মুকুল} (mukul) — a mango blossom; a bud}

\label{lesson:BN-C436-mukul}



\begin{quote}
Before the new one: say the Bengali for to trim, then the Bengali for a sprout.
\end{quote}

\subsection*{You'll want to know: \bn{মুকুল}}
\textbf{\bn{মুকুল}} — \emph{mukul} — \textquotedblleft{}a mango blossom; a bud\textquotedblright{}.

\begin{quote}
\bn{বসন্তে} \bn{আমগাছে} \bn{মুকুল} \bn{আসে।}
\end{quote}

\begin{quote}
\emph{bôsônte āmgāchhe mukul āse.}
\end{quote}

\begin{quote}
In spring the mango trees come into blossom.
\end{quote}

\begin{cousinweb}
\textbf{\bn{মুকুল}} (\emph{mukul}) is Sanskrit \emph{mukula}, \textquotedblleft{}a bud\textquotedblright{} (Monier-Williams); in Bengal it is above all the mango's spring blossom.
\end{cousinweb}

\begin{grammarlens}[title={What you've built}]
One more word for getting around.
\end{grammarlens}

\subsection*{Your turn}
\begin{itemize}
\raggedright
  \item \emph{Say it:} \emph{mukul}
  \item \emph{Say it:} \emph{mukul}, once more
  \item \emph{Say it:} \emph{bôsônte āmgāchhe mukul āse.} — In spring the mango trees come into blossom.
  \item \emph{Recall:} say the Bengali for to trim, then the Bengali for a sprout, then say \emph{mukul} again
\end{itemize}

\begin{tcolorbox}[breakable,colback=teal!4,colframe=teal!35!black,title={Before you move on}]
What does \bn{মুকুল} mean? (\textquotedblleft{}A mango blossom; a bud\textquotedblright{}.) Say it once more.
\end{tcolorbox}
\section[\texorpdfstring{\bn{ফোটা} (phoṭā)}{ফোটা (phoṭā)}]{\bn{ফোটা} (phoṭā) — to bloom; to open (of a flower); to boil}

\label{lesson:BN-C436-phota}



\begin{quote}
Before the new one: say the Bengali for a sprout, then the Bengali for a mango blossom.
\end{quote}

\subsection*{You'll want to know: \bn{ফোটা}}
\textbf{\bn{ফোটা}} — \emph{phoṭā} — \textquotedblleft{}to bloom; to open (of a flower); to boil\textquotedblright{}.

\begin{quote}
\bn{সকালে} \bn{বাগানে} \bn{গোলাপ} \bn{ফুটল।}
\end{quote}

\begin{quote}
\emph{sôkāle bāgāne golāp phuṭlô.}
\end{quote}

\begin{quote}
In the morning the roses bloomed in the garden.
\end{quote}

Past (he/she …): \bn{সেই} \bn{ফুল} \bn{ফুটল}, \emph{sei phul phuṭlô}.

The stem is \bn{ফুট}- in most forms: \bn{ফুটল}, \bn{ফুটবে}. Also of water coming to the boil: \bn{জল} \bn{ফুটল}.

\begin{cousinweb}
\textbf{\bn{ফোটা}} (\emph{phoṭā}) is inherited: CDIAL lists Bengali \emph{phuṭā}, \textquotedblleft{}to burst\textquotedblright{}, under the reconstructed \emph{sphuṭyati}, \textquotedblleft{}is burst open\textquotedblright{} (CDIAL 13845): a bud bursts open, and so does boiling water.
\end{cousinweb}

\begin{grammarlens}[title={What you've built}]
One more word for getting around.
\end{grammarlens}

\subsection*{Your turn}
\begin{itemize}
\raggedright
  \item \emph{Say it:} \emph{phoṭā}
  \item \emph{Say it:} \emph{phoṭā}, once more
  \item \emph{Say it:} \emph{sôkāle bāgāne golāp phuṭlô.} — In the morning the roses bloomed in the garden.
  \item \emph{Recall:} say the Bengali for a sprout, then the Bengali for a mango blossom, then say \emph{phoṭā} again
\end{itemize}

\begin{tcolorbox}[breakable,colback=teal!4,colframe=teal!35!black,title={Before you move on}]
What does \bn{ফোটা} mean? (\textquotedblleft{}To bloom; to open (of a flower); to boil\textquotedblright{}.) Say it once more.
\end{tcolorbox}
\section[\texorpdfstring{\bn{গজানো} (gôjāno)}{গজানো (gôjāno)}]{\bn{গজানো} (gôjāno) — to sprout; to grow (of leaves or hair)}

\label{lesson:BN-C436-gojano}



\begin{quote}
Before the new one: say the Bengali for a mango blossom, then the Bengali for to bloom.
\end{quote}

\subsection*{You'll want to know: \bn{গজানো}}
\textbf{\bn{গজানো}} — \emph{gôjāno} — \textquotedblleft{}to sprout; to grow (of leaves or hair)\textquotedblright{}.

\begin{quote}
\bn{বর্ষায়} \bn{উঠোনে} \bn{ঘাস} \bn{গজাল।}
\end{quote}

\begin{quote}
\emph{bôrshāy uṭhone ghās gôjālô.}
\end{quote}

\begin{quote}
In the monsoon grass sprang up in the courtyard.
\end{quote}

Past (he/she …): \bn{সেই} \bn{ঘাস} \bn{গজাল}, \emph{sei ghās gôjālô}.

\begin{cousinweb}
\textbf{\bn{গজানো}} (\emph{gôjāno}) is the everyday verb for leaves, grass or hair coming up; its origin is not settled.
\end{cousinweb}

\begin{grammarlens}[title={What you've built}]
One more word for getting around.
\end{grammarlens}

\subsection*{Your turn}
\begin{itemize}
\raggedright
  \item \emph{Say it:} \emph{gôjāno}
  \item \emph{Say it:} \emph{gôjāno}, once more
  \item \emph{Say it:} \emph{bôrshāy uṭhone ghās gôjālô.} — In the monsoon grass sprang up in the courtyard.
  \item \emph{Recall:} say the Bengali for a mango blossom, then the Bengali for to bloom, then say \emph{gôjāno} again
\end{itemize}

\begin{tcolorbox}[breakable,colback=teal!4,colframe=teal!35!black,title={Before you move on}]
What does \bn{গজানো} mean? (\textquotedblleft{}To sprout; to grow (of leaves or hair)\textquotedblright{}.) Say it once more.
\end{tcolorbox}
\section[\texorpdfstring{\bn{পোঁতা} (põtā)}{পোঁতা (põtā)}]{\bn{পোঁতা} (põtā) — to plant (in the ground); to bury}

\label{lesson:BN-C436-pota}



\begin{quote}
Before the new one: say the Bengali for to bloom, then the Bengali for to sprout.
\end{quote}

\subsection*{You'll want to know: \bn{পোঁতা}}
\textbf{\bn{পোঁতা}} — \emph{põtā} — \textquotedblleft{}to plant (in the ground); to bury\textquotedblright{}.

\begin{quote}
\bn{অরুণ} \bn{উঠোনে} \bn{একটা} \bn{আমগাছ} \bn{পুঁতল।}
\end{quote}

\begin{quote}
\emph{ôrun uṭhone ekṭā āmgāchh pũtlô.}
\end{quote}

\begin{quote}
Arun planted a mango tree in the courtyard.
\end{quote}

Past (he/she …): \bn{তিনি} \bn{গাছ} \bn{পুঁতলেন}, \emph{tini gāchh pũtlen}.

The stem is \bn{পুঁত}- in most forms: \bn{পুঁতি}, \bn{পুঁতল}; \bn{পোঁত}- in \bn{সে} \bn{পোঁতে}.

\begin{cousinweb}
\textbf{\bn{পোঁতা}} (\emph{põtā}) is inherited: CDIAL lists Bengali \emph{põtā}, \textquotedblleft{}to bury\textquotedblright{}, under the reconstructed \emph{pragupta}, \textquotedblleft{}hidden\textquotedblright{} (CDIAL 8474): to plant is to hide in the ground.
\end{cousinweb}

\begin{grammarlens}[title={What you've built}]
One more word for getting around.
\end{grammarlens}

\subsection*{Your turn}
\begin{itemize}
\raggedright
  \item \emph{Say it:} \emph{põtā}
  \item \emph{Say it:} \emph{põtā}, once more
  \item \emph{Say it:} \emph{ôrun uṭhone ekṭā āmgāchh pũtlô.} — Arun planted a mango tree in the courtyard.
  \item \emph{Recall:} say the Bengali for to bloom, then the Bengali for to sprout, then say \emph{põtā} again
\end{itemize}

\begin{tcolorbox}[breakable,colback=teal!4,colframe=teal!35!black,title={Before you move on}]
What does \bn{পোঁতা} mean? (\textquotedblleft{}To plant (in the ground); to bury\textquotedblright{}.) Say it once more.
\end{tcolorbox}
